\chapter{Contexte}

Au sein de l'ANTIC, les directions opérationnelles manipulent quotidiennement d'importants tableaux de bord sous forme de classeurs tableurs (\textbf{Microsoft Excel} ou \textbf{LibreOffice Calc}). Si le tableur est privilégié pour la manipulation des calculs et l'analyse numérique, la restitution officielle aux autorités exige la production de rapports formels rédigés sous \textbf{Microsoft Word}. Face aux limites du travail manuel actuel (recopie lourde, risques d'erreurs et désynchronisation des données), le présent projet s'inscrit dans la transformation numérique interne de l'institution pour automatiser le filtrage dynamique, la synchronisation en temps réel des feuilles enfants et la génération documentaire sécurisée.
